\documentclass[12pt]{article}
\usepackage{amsmath}
\usepackage{amssymb}
\usepackage{amsthm}
\usepackage{parskip}
\begin{document}
\noindent \textbf{Problem 1.2}
\smallbreak
\noindent \textbf{(a)}
\smallbreak
$\sqrt{(-1)(-1)} = \sqrt{(-1)}\sqrt{(-1)}$ is wrong, because: $$ a \ge 0 \lor b \ge 0 \implies \sqrt{a \cdot b} = \sqrt{a} \cdot \sqrt{b}$$
\noindent \textbf{(b)}
\smallbreak
\begin{proof}
Assume $1 = -1$ is true, we can add 3 to both sides of the equation yields: $$4 = 2$$ Then we divide both sides by 2, we obtain: $$2 = 1$$ which completes the proof.
\end{proof}

\noindent \textbf{(c)}
\smallbreak
\begin{proof}
    We want to show that $$\sqrt{r \cdot s} = \sqrt{r} \cdot \sqrt{s}$$
    Since $r$ and $s$ are positive real numbers as given, we can let $\sqrt{r} = a$ and $\sqrt{s} = b$, where $a$ and $b$ are positive real numbers.Squaring that we obtain $r = a^2$ and $s = b^2$, it follows that: $$r \cdot s = a^2 \cdot b^2 = (a \cdot b)^2$$
    Since $a$ and $b$ are positive real numbers, their product must also be a positive real number. Therefore, by the standard convention, $a \cdot b$ is the unique positive square root of $r \cdot s$, which means:$$\sqrt{r \cdot s} = a \cdot b$$
    Substituting $a = \sqrt{r}$ and $b = \sqrt{s}$ back into the equation, we obtain:$$\sqrt{r \cdot s} = \sqrt{r} \cdot \sqrt{s}$$
    which completes the proof.
\end{proof}
\end{document}

